% Source - https://tex.stackexchange.com/a/219399
% Posted by egreg, modified by community. See post 'Timeline' for change history
% Retrieved 2026-11-10, License - CC BY-SA 3.0

\documentclass[UTF8]{article}
\usepackage{amsmath}
\usepackage{amssymb}
\usepackage{graphicx}
\usepackage{epstopdf}
\usepackage{inputenc}
\usepackage{geometry}
\geometry{left=2.5cm,right=2.5cm,top=2.5cm,bottom=2.5cm}
\usepackage{hyperref}
\usepackage{amsfonts}

\begin{document}

\title{
  A Non-Regular Language in $O(\log \log n)$ Space \\
    \large
        Formalizing Analysis of Algorithms \\
        Project Proposal}

\author{Jora Zeneli, Laura Kribben, Lasse Reith, Elia Riva}

\date{13 November 2026}

\maketitle
\newpage

\section{Scope of the Project}

The project focuses on modeling a Turin Machine (TM), specifically a Multitape TM with two tapes: a read-only input tape and a work only input tape.
Both tapes may have different alphabets.
It is proven that Multitape TM reduce to single tape TM\cite{Intro-TOC-Sipser-1996} and it is easily seen that different alphabet for different tapes is reduced to a Multitape where in the working tape we simply never read or write symbols from the alphabet of the read-only tape.

The complexity measure we want to verify is the space complexity of a TM. Sipser\cite{Intro-TOC-Sipser-1996} defines the deterministic turning machine space complexity as: "Let $M$ be a deterministic Turing machine that halts on all inputs. The space complexity of $M$ is the function $f: \mathbb{N} \xrightarrow{} \mathbb{N},$ where $f(n)$ is the maximum number of tape cells that $M$ scans on any input of length $n$. If the space complexity of $M$ is $f(n)$, we also say that $M$ runs in space $f(n)$.

Once modeled we move into proving a series of theorems that will test the proposed modeling, finally proving the non-regular language requiring $O(\log \log n)$ space\cite{Lang-OneSymbAlph-LogLogn-HM-1975}

\section{Modeling Idea \& Objectives}

Modeling is the main task.

\section{Algorithms \& Complexity Results}

We will go over all algorithms that we will cover and state complexity results we wish to obtain.

\subsection{Accessory Algorithms}

Accessory algorithms plus expected results.
Accessory algorithms should be the simplest algorithm possible that go over edge cases of the model, so that we can verify exactly that we modeled the problem correctly.

\subsection{Preliminary Algorithms}

Preliminary algorithms plus expected results.
Preliminary algorithms should be medium to complex algorithms that bring some more 

\subsection{Main Algorithm}

Algorithm plus expected results

\section{Additional Tasks}

% TO IMPROVE and VERIFY
If we see that the scope of the project is not sufficient, we can model and verify that the language is not regular by defining a finite automata and verify that the language is not accepted, or go through the pumping lemma.

% max up to page 3
\newpage

% citations
\bibliographystyle{plain}
\bibliography{references}

\end{document}
